%\documentclass[12pt,a4paper,oneside]{report}
\documentclass[12pt,a4paper,oneside]{article} %se usiamo articole commentiamo \appendix
%\documentclass[a4paper, twoside]{article} %inizio documento, stampato su 2 lati, tipo articolo

\usepackage[italian]{babel} %imposta la lingua a italiano
\usepackage[utf8]{inputenc} %permette la compilazione delle lettere accentate

\usepackage{makeidx} %aggiunge il pakage degli indici
	\usepackage[breaklinks]{hyperref} %genera l'indice linkato
	\makeindex %genera l'indice

\usepackage{graphicx} %permette di importare i grafici in pdf e immagini
\usepackage{latexsym} %permette di inserire i caratteri speciali
\usepackage{multirow} %colonne a multiriga
\usepackage{verbatim} %commento a blocco usando \begin{comment} \end{comment}

\usepackage[toc,acronym]{glossaries}
	% --- ACRONIMI ---

\newacronym{cas}{CAS}{\gls{Central Authentication Service}}
\newacronym{dms}{DMS}{\gls{Document Management System}}
\newacronym{crud}{CRUD}{Create Read Update Delete}

% --- GLOSSARIO ---

\newglossaryentry{Central Authentication Service}
{
  name=Central Authentication Service,
  description={è un servizio single sign-on libero che permette alle applicazioni web la possibilità di rinviare tutte le autenticazioni a un server centrale o a più server di fiducia}
}
\newglossaryentry{Document Management System}
{
  name=Document Management System,
  description={è una categoria di sistemi software che serve a organizzare e facilitare la creazione collaborativa di documenti e di altri contenuti.\\
Tecnicamente il DMS è un'applicazione lato server che si occupa di eseguire operazioni massive sui documenti, catalogandoli ed indicizzandoli secondo determinati algoritmi}
}
	\makeglossaries

\pagestyle{plain} %stampa indice a piepagina

\title{Moddlepedia - Wikiddle - }

\author{Amodio\\Stefano\\727457
\and Pinzan\\Andrea\\732345
\and Racchetti\\Luca\\703311}

\date{\today}

\begin{document} %inizio del documento

%-------------------------------------------------------------------------

\begin{titlepage}
%\maketitle %genera la prima pagina
\begin{center}
		\Huge{Moddlepedia - Wikiddle}\\
	\vskip 20pt
	\begin{Large}
		\begin{tabular}{ c c c }
		  Amodio & Pinzan & Racchetti \\
		  Stefano & Andrea & Luca \\
		  727457 & 732345 & 703311 \\
		\end{tabular}
	\end{Large}\\
	\vskip 20pt
		\begin{normalsize}
			\today\\
		\end{normalsize}
	\vfill
		\includegraphics[scale=0.8]{Immagini/logo.jpg}
	\vfill
	\begin{normalsize}
		Versione 0.0.1	
	\end{normalsize}
\end{center}
\end{titlepage}

\thispagestyle{empty}
\newpage %nuova pagina

\tableofcontents
\appendix
\printindex % stampa l'indice
\newpage %nuova pagina
%-------------------------------------------------------------------------



%-------------------------------------------------------------------------

\section{Ambito e Scopo del documento e del progetto}

\gls{cas} è un acronimo la prima volta.\\
\gls{cas} è un acronimo la seconda volta.\\
\gls{dms}\\
\gls{Central Authentication Service} è un glossario la prima volta.\\
\gls{Central Authentication Service} è un glossario la seconda volta.

aaaaa \cite{esp1} bbbb \cite{esp1} cccc

\newpage %nuova pagina

%-------------------------------------------------------------------------
	
\section{Acronimi e Glossario}

\deftranslation{Glossary}{Glossario}
\deftranslation{Acronyms}{Acronimi}
\renewcommand*{\glossaryname}{Glossario}
\renewcommand*{\acronymname}{Acronimi}

\setglossarysection{subsection}

%\glsaddall %imposta la stampa di tutti i termini del gossario, anche se non utilizzati

%\addcontentsline{acronym}{subsection}{Acronimi}
%\printglossary[type=acronym]
%\printglossary[type=acronym,title=Acronimi]

%\addcontentsline{toc}{subsection}{Glossario}
%\printglossary[type=toc]
%\printglossary[type=toc,title=Glossario]

\printglossary[numberedsection,type=acronym]
\printglossary[numberedsection]
%\printglossaries

\newpage %nuova pagina

%-------------------------------------------------------------------------
\section{Documenti utilizzati}

\newpage %nuova pagina
%-------------------------------------------------------------------------

\section{Gap Analysis}

\newpage %nuova pagina
%-------------------------------------------------------------------------

\section{Modello degli oggetti organizzativi (Business Object Model)}

\subsection{Analisi di Dominio - Slide 13}

\newpage %nuova pagina	

%-------------------------------------------------------------------------

\section{Elenco dei requisiti funzionali e informativi}

\subsection{Req.1 - Gestione documento}
\paragraph{Tipologia:}
	(funzionale, non funzionale, informativo...)
\paragraph{Descrizione:}
	L'utente deve poter effettuare il \gls{crud} su un documento.
\paragraph{Importanza:}
	N (Need)
\paragraph{Motivazione:}
	Le azioni di \gls{crud} su un documento sono necessarie per poter poterlo gestire.
\paragraph{Criterio di Validazione:}
	Il documento può essere creato?\\
	Il documento può essere modificato?\\
	Il documento può essere visualizzato?\\
	Il documento può essere cancellato?
\paragraph{Legame con altri requisiti:}
	La gestione di un documento è la base per requisiti successivi.

\subsection{Req.2 - Caricamento file utilizzabile nel documento}
\paragraph{Tipologia:}
	(funzionale, non funzionale, informativo...)
\paragraph{Descrizione:}
	L'utente deve poter caricare un file (immagine, pdf, audio, video, ...) utilizzabile nel corpo del documento.
\paragraph{Importanza:}
	N (Need)
\paragraph{Motivazione:}
	Rendere un file disponibile all'interno del documento a fine esplicativo.
\paragraph{Criterio di Validazione:}
	Il file può essere caricato?\\
	Il file è disponibile nel documento?
\paragraph{Legame con altri requisiti:}
	\subparagraph{Prerequisito} Req.1


\subsection{Req.3 - Gestione di gruppi di utenti}
\paragraph{Tipologia:}
	(funzionale, non funzionale, informativo...)
\paragraph{Descrizione:}
	Il docente può creare gruppi di utenti


\subsection{Req.4 - Gestione dei permessi sul documento}
\subsection{Req.5 - Gestione dei permessi sul file}







\subsection{Req.6 - Gestione delle versioni di un documento}
\paragraph{Tipologia:}
	(funzionale, non funzionale, informativo...)
\paragraph{Descrizione:}
	L'utente deve poter inserire un file (immagine, pdf, audio, video, ...) nel corpo del documento.
\paragraph{Importanza:}
	N (Need)
\paragraph{Motivazione:}
	Rendere un file disponibile all'interno del documento a fine esplicativo.
\paragraph{Criterio di Validazione:}
	Il file può essere caricato?\\
	Il file viene inserito?\\
	Il file è collocato dove indicato?
\paragraph{Legame con altri requisiti:}
	\subparagraph{Prerequisito} Req.1 \\!!TO-DO

\subsection{Req.7 - Gestione delle versioni di un file}





\newpage %nuova pagina

%-------------------------------------------------------------------------	

\section{Casi d'uso e di interazione}

\subsection{CASO1}
\subsubsection{CASO1 - Obiettivo}
\subsubsection{CASO1 - Descrizione ordinata e dettagliata interazione fra Utenti o Utenti-Sistema}
\subsubsection{CASO1 - Scenari}
\paragraph{CASO1 - Happy Flow}
	\begin{itemize}
	\item acquisto iniziato e concluso con successo
	\end{itemize}
\paragraph{CASO1 - Percorsi alternativi}
	\begin{itemize}
	\item Disponiblità di un prodotto terminata
	\item Carta di credito non acettata
	\item Collegamento con servizi interbancari interrotto
	\end{itemize}
\subsubsection{CASO1 - Diagrammi}

\newpage %nuova pagina

%-------------------------------------------------------------------------

\section{Descrizione di un processo notevole di interazione}

\subsection{Modello/diagramma dinamico (processo e di collaborazione)}

\newpage %nuova pagina

%-------------------------------------------------------------------------
%-------------------------------------------------------------------------
%-------------------------------------------------------------------------
%-------------------------------------------------------------------------
%-------------------------------------------------------------------------
%-------------------------------------------------------------------------
%-------------------------------------------------------------------------
%-------------------------------------------------------------------------

\section{Diagramma}

\subsection{di Struttura}

\subsubsection{static structure diagram}
\subsubsection{entity diagram}
\subsubsection{class diagram}
\subsubsection{activity diagram / behavior}
\subsubsection{BPMN}

\newpage %nuova pagina

%-------------------------------------------------------------------------
	
\section{Appendici}

\subsection{diagrammi di sequenza}
\subsection{collaboration diagram}
\subsubsection{entity diagram + scambi di messaggio}

\newpage %nuova pagina

%-------------------------------------------------------------------------


\begin{thebibliography}{50}
 \bibitem{esp1} D. P. Woodruff, W. A. Royer, N. V. Smith (1986), \emph{Physical Review B} {\bf 34}, 2.
 \bibitem{esp2} K. Giesen, F. Hage, F. J. Himpsel, H. J. Riess, W. Steinmann (1986), \emph{Physical Review B} {\bf 33}, 8.
  \bibitem{esp3} M. Ortuno, P. M. Echenique (1986), \emph{Physical Review B} {\bf 34}, 8.
  \bibitem{A-B} J. A. Appelbaum, E. I. Blount (1973), \emph{Physical Review B} {\bf 8}, 483.
 \bibitem{Art:PG} J. B. ~Pendry, S. J. Gurman (1974), \emph{Surface Science } {\bf 49}, 87-105.   
 \bibitem{Art:Griffiths} David J.~Griffiths (2005), \emph{Introduzione alla Meccanica Quantistica}, Ambrosiana.
\bibitem{Acciarri} N. W. Ashcroft, N. D. Mermin (1976), \emph{Solid State Physics}, Holt, Rinehart, Winston.
	\bibitem {2bande} P. M. Echenique, J. M. Pitarke, E. V. Chulkov, V. M. Silkin (2008), \emph{Journal of electron spectroscopy and related phenomena} {\bf 126}, 163-175.
\bibitem{Art:Chulkov} E. V. ~Chulkov, V. M. Silkin, P. M. Echenique (1999), \emph{Surface Science} {\bf 437}, 330-352.
 \bibitem{Cohen:quantum1} C.~Cohen-Tannoudji, B.~Diu, F.Laloe (1973), \emph{Quantum Mechanics}, Vol.~1, Wiley-VCH.
\bibitem{bucatri}  J. H. Davies (1998), \emph{The Physiscs of Low Dimensional Semiconductors-An Introduction}, Cambridge University Press.
   \bibitem{Cohen:quantum2} C.~Cohen-Tannoudji, B.~Diu, F.Laloe (1973), \emph{Quantum Mechanics}, Vol.~2, Wiley-VCH.
 \bibitem{Art:EP} P. M. ~Echenique, J. B. Pendry (1977), \emph{Solid State Physics} {\bf 11}, 2065.
 \bibitem{Art:Smith} N. V. ~Smith (1985), \emph{Physical Review B} {\bf 32}, 6.
   \bibitem {} G.Butti (2005) \emph{PhD Thesis: Surface specific electronic structure on extended substrates}.

\end{thebibliography}


\end{document} %fine del documento
